\renewcommand{\theposttitle}{How we frame failure paves our road to success}
\renewcommand{\thepostdate}{2024-01-16}
\vspace*{\fill}
\begin{center}
\includegraphics[width=\textwidth, height=0.55\textheight, keepaspectratio]{C:/Users/James/Documents/GitHub/VisualIdeasBook/data/images/28f2ba379c0f25f5b675e15ffa6db341.png}
\end{center}
\vspace{2mm}
\begin{center}
\begin{minipage}{0.92\textwidth}
\small
The mental model that stops me from having an opinion about everything:

\textbf{Dunning-Kruger Effect}--\\
The less you know, the less you question.\\
The more you know, the more you question everything.

I’m fascinated by self-deception and our lack of awareness of how we think and make decisions. And the Dunning-Kruger Effect is a self-deception we all can fall prey to…because well, you don’t know what you don’t know.

One of my goals this year is to reread (and take notes this time!) the book, \textbf{“Thinking Fast and Thinking Slow” by Daniel Kahneman, }because it’s mentioned in practically every book on psychology and sociology ever written since it’s inception. Watch this space for more visuals on cognitive biases 👀
\end{minipage}
\end{center}
\vspace*{\fill}